\section{Cuerpos arquimedianos}

Si $\K$ es un cuerpo ordenado, se tiene $\N \subset \Z \subset \Q \subset \K$, pero esto no implica que $\N$ sea no acotado en $\K$.

\begin{eg}
  Sea $\Q(t)$ el cuerpo de funciones racionales $r(t) = p(t)/q(t)$, con $p, q \in \Q[t]$ y $q \neq 0$. Se lo ordena declarando que $r(t)$ es positivo si y solo si el coeficiente principal del polinomio $p \cdot q$ es positivo.

  Para cada $n \in \N$, el polinomio $t - n$ tiene coeficiente principal $1 > 0$, de modo que $t - n \in P$, es decir, $t > n$. Por lo tanto, en $\Q(t)$ el conjunto $\N$ está acotado superiormente, por ejemplo por $t$.
\end{eg}

\begin{definition}[Cuerpo arquimediano]
  Un cuerpo ordenado $\K$ es \emph{arquimediano} si $\N \subset \K$ no es acotado superiormente.
\end{definition}

\begin{theorem}
  Sea $\K$ un cuerpo ordenado. Son equivalentes:
  \begin{enumerate}[label=(\arabic*)]
    \item $\K$ es arquimediano, es decir, $\N$ no es acotado superiormente.
    \item Dados $a, b \in \K$ con $a > 0$, existe $n \in \N$ tal que $n a > b$.
    \item Dado $a \in \K$ con $a > 0$, existe $n \in \N$ tal que $0 < \frac{1}{n} < a$.
  \end{enumerate}
\end{theorem}

\begin{proof}
  \begin{description}[font=\normalfont, style=nextline, leftmargin=2em, itemsep=0.3em]
    \item[(1) $\Rightarrow$ (2)] Como $\N$ no es acotado, $b/a$ no es cota superior de $\N$, así que existe $n \in \N$ con $n > b/a$. Luego $na > b$.
    \item[(2) $\Rightarrow$ (3)] Dado $a > 0$, aplicamos (2) con $b = 1$: existe $n$ con $na > 1$, es decir, $0 < \frac{1}{n} < a$.
    \item[(3) $\Rightarrow$ (1)] Sea $b \in \K$. Si $b \leq 0$, claramente $b$ no es cota superior de $\N$. Si $b > 0$, aplicamos (3) con $a = 1/b$: existe $n$ con $0 < \frac{1}{n} < \frac{1}{b}$, o sea $n > b$. Luego ningún elemento de $\K$ es cota superior de $\N$.\qedhere
  \end{description}
\end{proof}

\section{Supremo e ínfimo}

\begin{definition}[Supremo]
  Sean $\K$ un cuerpo ordenado y $X \subset \K$ acotado superiormente. Un elemento $b \in \K$ es el \emph{supremo} de $X$ si es la menor de las cotas superiores de $X$ en $\K$. Es decir:
  \begin{enumerate}[label=(\arabic*)]
    \item $x \leq b$ para todo $x \in X$;
    \item si $c \in \K$ y $x \leq c$ para todo $x \in X$, entonces $b \leq c$.
  \end{enumerate}
  Equivalentemente, (2) puede reemplazarse por: para todo $c < b$ existe $x \in X$ tal que $c < x$.
\end{definition}

\begin{note}
  \begin{enumerate}[label=(\arabic*)]
    \item El supremo de un conjunto, si existe, es único.
    \item Si un conjunto tiene máximo, este es su supremo.
    \item Si $X = \varnothing$, todo elemento de $\K$ es cota superior. Como $\K$ no tiene mínimo, $\varnothing$ no tiene supremo en $\K$.
  \end{enumerate}
\end{note}

\begin{definition}[Ínfimo]
  Sean $\K$ un cuerpo ordenado y $X \subset \K$ acotado inferiormente. Un elemento $b \in \K$ es el \emph{ínfimo} de $X$ si es la mayor de las cotas inferiores de $X$ en $\K$. Es decir:
  \begin{enumerate}[label=(\arabic*)]
    \item $x \geq b$ para todo $x \in X$;
    \item si $c \in \K$ y $x \geq c$ para todo $x \in X$, entonces $b \geq c$.
  \end{enumerate}
  Equivalentemente, (2) puede reemplazarse por: para todo $c > b$ existe $x \in X$ tal que $x < c$.
\end{definition}

\begin{eg}
  Dados $a < b$ en $\K$, si $X = (a, b)$ entonces $\inf X = a$ y $\sup X = b$.

  Por definición, $a$ es cota inferior y $b$ es cota superior. Veamos que $a$ es la mayor cota inferior. Sea $c > a$. Si $c \geq b$, entonces $c$ no es cota inferior pues cualquier $x \in X$ cumple $x < b \leq c$. Si $c < b$, entonces $\frac{a + c}{2} \in X$ y $\frac{a + c}{2} < c$, de modo que $c$ tampoco es cota inferior. Luego $a = \inf X$. El argumento para el supremo es análogo.
\end{eg}

\begin{eg}
  Sea $Y = \left\{ \frac{1}{2^n} : n \in \N \right\} \subset \Q$. Entonces $\sup Y = \frac{1}{2}$ y $\inf Y = 0$.

  Como $\frac{1}{2} \in Y$ y $\frac{1}{2^n} \leq \frac{1}{2}$ para todo $n$, se tiene $\sup Y = \frac{1}{2}$.

  Además $0 < \frac{1}{2^n}$ para todo $n$, así que $0$ es cota inferior. Sea $c > 0$. Por la desigualdad de Bernoulli, $2^n = {(1 + 1)}^n \geq 1 + n > n$. Como $\Q$ es arquimediano, existe $n$ con $n > \frac{1}{c}$, y entonces $2^n > \frac{1}{c}$, es decir, $\frac{1}{2^n} < c$. Luego $c$ no es cota inferior y $\inf Y = 0$.
\end{eg}

El problema más serio de los racionales desde el punto de vista del análisis es que algunos conjuntos acotados de números racionales no tienen supremo (o ínfimo) en $\Q$.

\begin{eg}
  Sean
  \[
    X = \{ x \in \Q : x \geq 0,\ x^2 < 2 \}, \qquad Y = \{ y \in \Q : y > 0,\ y^2 > 2 \}.
  \]
  Si $z > 2$, entonces $z^2 > 4$ y $z \notin X$; luego $X \subset [0, 2]$ está acotado. Además $Y \subset (0, +\infty)$ está acotado inferiormente. Veamos que $X$ no tiene supremo en $\Q$.

  \begin{enumerate}[label=(\arabic*)]
    \item \emph{$X$ no tiene máximo.} Dado $x \in X$ se elige $r \in \Q$ con $0 < r < \min\left\{1, \frac{2 - x^2}{2x + 1}\right\}$. Como $r < 1$, resulta $r^2 < r$ y
          \[
            {(x + r)}^2 = x^2 + 2xr + r^2 < x^2 + r(2x + 1) < 2,
          \]
          de modo que $x + r \in X$ y $x + r > x$.
    \item \emph{$Y$ no tiene mínimo.} Dado $y \in Y$ se elige $r \in \Q$ con $0 < r < \frac{y^2 - 2}{2y}$. Notar que $\frac{y^2-2}{2y} < \frac{y}{2}$, de modo que $y - r > 0$. Además
          \[
            {(y - r)}^2 = y^2 - 2yr + r^2 > y^2 - 2yr > y^2 - (y^2 - 2) = 2,
          \]
          así que $y - r \in Y$ e $y - r < y$.
    \item \emph{Si $x \in X$ e $y \in Y$, entonces $x < y$.} En efecto, si fuese $y \leq x$, sería $y^2 \leq x^2 < 2$, lo que contradice $y^2 > 2$.
  \end{enumerate}

  Supongamos que existe $\alpha = \sup X \in \Q$. Como $1 \in X$, es $\alpha \geq 1 > 0$.
  \begin{itemize}
    \item Si $\alpha^2 < 2$, entonces $\alpha \in X$ sería el máximo de $X$, lo que contradice (1).
    \item Si $\alpha^2 > 2$, entonces $\alpha \in Y$. Por (2) existe $\beta \in Y$ con $\beta < \alpha$ y por (3) $\beta$ es cota superior de $X$, lo que contradice que $\alpha$ sea la menor cota superior.
  \end{itemize}
  Por lo tanto $\alpha^2 = 2$. Pero esto es imposible en $\Q$: si ${(p/q)}^2 = 2$ con $p, q$ coprimos, entonces $p^2 = 2q^2$ es par, luego $p$ es par, y entonces $4 \mid 2q^2$ implica que $q$ también es par, absurdo. Así $X$ no tiene supremo en $\Q$. El caso de $\inf Y$ es análogo.
\end{eg}

\section{Cuerpos completos}

Si $\K$ es un cuerpo ordenado no arquimediano, $\N$ está acotado superiormente. Pero además no tiene supremo: si $b$ es una cota superior de $\N$, entonces $n + 1 \leq b$ para todo $n \in \N$, es decir, $n \leq b - 1$ para todo $n$, de modo que $b - 1 < b$ también es cota superior.

\begin{definition}[Cuerpo completo]
  Un cuerpo ordenado $\K$ es \emph{completo} si todo subconjunto no vacío y acotado superiormente de $\K$ tiene supremo en $\K$.
\end{definition}

\begin{prop}
  Sea $\K$ un cuerpo ordenado completo. Entonces:
  \begin{enumerate}[label=(\arabic*)]
    \item todo subconjunto no vacío y acotado inferiormente tiene ínfimo en $\K$;
    \item $\K$ es arquimediano.
  \end{enumerate}
\end{prop}

\begin{proof}
  (1) Sea $Y \subset \K$ no vacío y acotado inferiormente. Entonces $X = -Y = \{ -y : y \in Y \}$ es no vacío y acotado superiormente, de modo que existe $a = \sup X$, y se verifica que $-a = \inf Y$.

  (2) Si $\N$ estuviese acotado superiormente, tendría un supremo $s$. Como $s - 1 < s$ no es cota superior de $\N$, existe $n \in \N$ con $n > s - 1$, y entonces $n + 1 > s$ con $n + 1 \in \N$, lo que contradice que $s$ sea cota superior.
\end{proof}

\begin{note}
  Se asume como axioma que existe un cuerpo ordenado completo, que llamamos $\R$ (y que es único salvo isomorfismo).
\end{note}

\begin{ejercicio}
  Dados $a \in \R$ con $a > 0$ y $n \in \N$, existe un único $b > 0$ tal que $b^n = a$.

  \emph{Sugerencia.} Definir $X = \{ x \in \R : x \geq 0,\ x^n < a \}$ e $Y = \{ y \in \R : y > 0,\ y^n > a \}$ e imitar el ejemplo anterior. Para ello, probar y usar que, dado $x > 0$, existe una constante $A_n(x) > 0$ tal que ${(x + d)}^n \leq x^n + A_n(x)\, d$ para todo $0 < d < 1$.
\end{ejercicio}

\begin{definition}[Conjunto denso]
  Un conjunto $X \subset \R$ es \emph{denso} en $\R$ si todo intervalo abierto $(a, b)$ con $a < b$ contiene algún punto de $X$.
\end{definition}

\begin{prop}
  $\Q$ es denso en $\R$.
\end{prop}

\begin{proof}
  Sean $a < b$. Como $b - a > 0$ y $\R$ es arquimediano, existe $p \in \N$ con $0 < \frac{1}{p} < b - a$. Sea $A = \{ m \in \Z : \frac{m}{p} \geq b \}$. Por la propiedad arquimediana $A \neq \varnothing$, y está acotado inferiormente por $bp$. Entonces $A$ tiene un mínimo $m_0$, que cumple $\frac{m_0}{p} \geq b$ y $\frac{m_0 - 1}{p} < b$.

  Veamos que $\frac{m_0 - 1}{p} > a$. Si no, tendríamos $\frac{m_0 - 1}{p} \leq a < b \leq \frac{m_0}{p}$, y entonces $b - a \leq \frac{m_0}{p} - \frac{m_0 - 1}{p} = \frac{1}{p}$, lo que es absurdo. Luego $\frac{m_0 - 1}{p} \in (a, b) \cap \Q$.
\end{proof}

\begin{prop}
  $\R \setminus \Q$ es denso en $\R$.
\end{prop}

\begin{proof}
  Sean $a < b$. Como $\Q$ es denso, existe $r \in \Q$ con $a - \sqrt{2} < r < b - \sqrt{2}$. Entonces $r + \sqrt{2} \in (a, b)$ y $r + \sqrt{2} \notin \Q$, pues de lo contrario $\sqrt{2} = (r + \sqrt{2}) - r$ sería racional.
\end{proof}

\section{Introducción a la cardinalidad}

\begin{definition}[Equipotencia]
  Dos conjuntos $X$ e $Y$ tienen el \emph{mismo cardinal} (o son \emph{coordinables} o \emph{equipotentes}) si existe una función biyectiva $f\colon X \to Y$. Se escribe $X \sim Y$ o $\card(X) = \card(Y)$.
\end{definition}

La relación $\sim$ es de equivalencia.

Dado $n \in \N$, sea $I_n = \{1, 2, \dots, n\}$.

\begin{lemma}
  Para todo $n \in \N$ no existe una función suryectiva $h\colon I_n \to I_{n+1}$.
\end{lemma}

\begin{proof}
  Por inducción en $n$. Para $n = 1$, $I_1$ tiene un solo elemento, así que su imagen no puede ser $I_2$.

  Supongamos que el resultado vale para $n - 1$ y que existe $h\colon I_n \to I_{n+1}$ suryectiva. Sea $r$ la permutación de $I_{n+1}$ que intercambia $h(n)$ con $n + 1$. Entonces $r \circ h$ es suryectiva y $(r \circ h)(n) = n + 1$. Cada $j \in I_n$ es imagen de algún $k \in I_n$ y necesariamente $k \neq n$ (pues $(r \circ h)(n) = n + 1 \notin I_n$). Por lo tanto, la restricción $(r \circ h)|_{I_{n-1}}\colon I_{n-1} \to I_n$ es suryectiva, lo que contradice la hipótesis inductiva.
\end{proof}

\begin{theorem}
  Sean $n, m \in \N$. Entonces $I_n \sim I_m$ si y solo si $n = m$.
\end{theorem}

\begin{proof}
  Si $n = m$ es trivial. Recíprocamente, sea $f\colon I_n \to I_m$ biyectiva y supongamos, sin pérdida de generalidad, que $n < m$ (si no, se usa $f^{-1}$). Definimos $g\colon I_m \to I_{n+1}$ por
  \[
    g(k) =
    \begin{cases}
      k & \text{si } 1 \leq k \leq n + 1, \\
      1 & \text{si } k > n + 1.
    \end{cases}
  \]
  Esta función es suryectiva y, como $f$ es biyectiva, $g \circ f\colon I_n \to I_{n+1}$ también lo es, lo que contradice el lema.
\end{proof}

\begin{definition}[Conjunto finito]
  Un conjunto $X$ es \emph{finito} si es vacío o si existe $n \in \N$ tal que $X \sim I_n$; en ese caso $\card(X) = n$ (y $\card(\varnothing) = 0$). Por el teorema anterior, $n$ está bien definido. Un conjunto que no es finito se dice \emph{infinito}.
\end{definition}
